% Part: incompleteness
% Chapter: introduction
% Section: overview

\documentclass[../../../include/open-logic-section]{subfiles}

\begin{document}

\olfileid{inc}{int}{ovr}

\olsection{د ناتکميلۍ د پايلو ټوليزه کتنه}

هېلبرټ تمه لرله چې رياضيات په داسې !!{axiomatizable} تيورۍ کښې
صوري کېدے شي چې !!{complete} او !!{decidable} ثابتېدے شي۔ سربېره پر
دې، موخه ئې دا وه چې د دې تيورۍ سازګاري په ډېرو کمزورو،
«متناهي‌پالو» لارو ثابته کړي، څو د شهوديت د ننګونو پر وړاندې له
کلاسیکو رياضياتو دفاع وکړي۔ د ګوډل د ناتکميلۍ قضيو وښودل چې دا موخې
نۀ شي تر لاسه کېدے۔

د ګوډل د ناتکميلۍ لومړۍ قضيې وښودل چې د رسل او وايټ‌هېډ د
\emph{پرېنسيپيا مېتېماتيکا} يوه بڼه !!{complete} نۀ ده۔ خو ثبوت په
اصل کښې ډېر عام و او پر ډېرو بېلابېلو تيوريو پلي کېږي۔ معنا ئې دا ده
چې ستونزه يوازې دا نۀ وه چې \emph{پرېنسيپيا مېتېماتيکا} رياضيات په
بشپړ ډول رانۀ غښتل، بلکې \emph{هېڅ} منل کېدونکې تيوري ئې نۀ شي
رانغاړلے۔ څه وخت ولګېد چې د تيوريو هغه ځانګړتياوې بېلې شي چې د
ناتکميلۍ د قضيو د پلي کېدو دپاره بس دي، او د ګوډل ثبوت داسې عام شي
چې يوازې په همدې ځانګړتياوو تکيه وکړي۔ خو اوس کولے شو د حساب د
ژبې~$\Lang L_A$ د تيوريو دپاره د ناتکميلۍ د لومړۍ قضيې يوه ډېره عامه
بڼه بيان کړو۔

\begin{thm}
که~$\Gamma$ په~$\Lang L_A$ کښې داسې سازګاره او !!{axiomatizable}
تيوري وي چې ټولې محاسبه کېدونکې تابعې او ټولې !!{decidable} اړيکې
!!{represents} کوي، نو~$\Gamma$ !!{complete} نۀ ده۔
\end{thm}

دا ويل چې~$\Gamma$ !!{complete} نۀ ده، په دې معنا دي چې لږ تر لږه
د يوې !!{sentence}~$!A$ دپاره~$\Gamma \Proves/ !A$ او
$\Gamma \Proves/ \lnot !A$ وي۔ داسې !!a{sentence} (له~$\Gamma$
څخه) \emph{خپلواکه} بلل کېږي۔
\textbf{د ژباړې د نسخې سمون (OLCMP-070):} سرچينې د ګاما له نوم
سره د تړون قوس په ناسمه توګه د رياضياتي چاپېريال دننه ليکلے و؛ دلته
قوس د جملې په خپل ځای کښې راوړل شوے او د تيورۍ نوم نۀ دے بدل شوے۔
په حقيقت کښې په نسبتاً چټکۍ ثابتولے شو
چې خپلواکې جملې بايد موجودې وي۔ خو د قضيې د ګوډلي ثبوت ځواک په دې کښې
دے چې د داسې خپلواکې !!{sentence} يوه \emph{ټاکلې بېلګه} وړاندې کوي۔
دا په زړه پورې جوړونه يوه !!a{sentence}~$!G_\Gamma$ پيدا کوي، چې د
$\Gamma$ \emph{ګوډل جمله} بلل کېږي او ځکه ناثابتېدونکې ده چې په
$\Gamma$ کښې~$!G_\Gamma$ له دې ادعا سره برابره ده چې~$!G_\Gamma$ په
$\Gamma$ کښې ناثابتېدونکې ده۔ جوړونه دا کار \emph{تعميري} کوي؛ يعنې
که د~$\Gamma$ بديهي کول او د !!{derivation} نظام بيان راکړل شي، ثبوت
په رښتيا د~$!G_\Gamma$ د ليکلو يوه لار راکوي۔

د ګوډل په ثبوت کښې جوړونه غواړي چې د~$\Lang L_A$ د اصطلاحاتو او
!!{formula}s ځانګړتياوې او عملونه پخپله په~$\Lang L_A$ کښې د بيانولو
لار پيدا کړو۔ په دې کښې داسې ځانګړتياوې شاملې دي لکه «$!A$
!!a{sentence} ده»، «$\delta$ د~$!A$ !!a{derivation} ده»، او داسې
عملونه لکه~$\Subst{!A}{t}{x}$۔ دا لار بايد (الف) دا ځانګړتياوې او
اړيکې د نښو او د هغو د لړيو د «کوډولو» په وسيله د طبيعي عددونو په
توګه بيان کړي (ځکه اصطلاحات، !!{formula}s، !!{derivation}s او نور د
نښو لړۍ دي، او~$\Lang L_A$ د طبيعي عددونو په باب خبرې کوي)۔ بايد
(ب) دا کار داسې وکړي چې~$\Gamma$ اړوند حقايق ثابت کړي؛ نو بايد وښيو
چې دا ځانګړتياوې د طبيعي عددونو په !!{decidable} خاصيتونو کوډېږي او
عملونه د طبيعي عددونو له محاسبه کېدونکو تابعو سره سمون لري۔ دې ته
«د نحو حسابي کول» ويل کېږي۔

خو د نحو د حسابي کولو تر څېړلو وړاندې به دا شرط وګورو چې~$\Gamma$
«کافي پياوړې» وي؛ يعنې ټولې محاسبه کېدونکې تابعې او ټولې
!!{decidable} اړيکې !!{represents} کړي۔ دا غواړي چې «محاسبه کېدونکي»
ته يو کره تعريف ورکړو۔ دا په بېلابېلو لارو کېدے شي؛ مثلاً د ټيورينګ
ماشينونو د مدل په وسيله، يا د هغو تابعو په توګه چې په کومې عمومي
موخې پروګرام‌ليک ژبه کښې پروګرامونه ئې محاسبه کوي۔ خو څرنګه چې زموږ
موخه د~$\Lang L_A$ د ژبې په يوه تيورۍ کښې د دې تابعو او اړيکو
!!{represents} کول دي، غوره ده چې يوازې د عددي تابعو د محاسبه کېدنې يو ساده تعريف
وټاکو۔ دا د \emph{بازګشتي تابعې} مفهوم دے۔ نو لومړے به د بازګشتي
تابعو بحث وکړو۔ بيا به وښيو چې~$\Th{Q}$ لا له وړاندې ټولې بازګشتي
تابعې او اړيکې !!{represents} کوي۔ دا به موږ ته اجازه راکړي چې د
ناتکميلۍ قضيه پر~$\Th{Q}$ او~$\Th{PA}$ غوندې ټاکلو تيوريو پلې کړو،
ځکه ثابته کړې به مو وي چې دا د «کافي پياوړو» تيوريو بېلګې دي۔
\textbf{د ژباړې د نسخې سمون (OLCMP-071):} سرچينې د دې فقرې په
موخه‌بيانوونکې جمله کښې د اصطلاحي نښې په پای کښې زائده انګرېزي~s
لرله؛ دلته هماغه ټاکل شوې اصطلاح له زائدې نښې پرته راوړل شوې ده۔

د نحو د حسابي کولو وروستۍ پايله داسې !!a{formula}~$\OProv[\Gamma](x)$
ده چې د !!{formula}s د عددونو په توګه د کوډولو په وسيله د~$\Gamma$
له بديهي اصولو څخه د اثبات وړتيا بيانوي۔ په ځانګړي ډول، که~$!A$ په
عدد~$n$ کوډ شوې وي او~$\Gamma \Proves !A$، نو
$\Gamma \Proves \Prov[\Gamma](\num{n})$۔ د~$\Gamma$ دپاره دا
«اثبات‌وړتيا پريديکات» موږ ته دا وس هم راکوي چې په يوه معنا د~$\Gamma$
سازګاري د~$\Lang L_A$ د يوې !!a{sentence} په توګه بيان کړو: د
$\Gamma$ دپاره د «سازګارۍ جمله» په نوم هغه !!{sentence}
$\lnot \Prov[\Gamma](\num{n})$ ونيسئ، چې~$n$ پکې د يوه تناقض،
مثلاً~$\lfalse$، کوډ وي۔ د ناتکميلۍ
دويمه قضيه وايي چې سازګارې !!{axiomatizable} تيورۍ خپلې د سازګارۍ
جملې هم نۀ ثابتوي۔ د دې قضيې د پلي کېدو شرطونه تر دې لږ سخت دي چې
تيوري دې يوازې ټولې محاسبه کېدونکې تابعې او !!{decidable} اړيکې
تمثيل کړي، خو وبه ښيو چې~$\Th{PA}$ دا شرطونه پوره کوي۔

\end{document}
